% Generated draft: requires mathematical and source review.
\documentclass[11pt]{article}
\usepackage[margin=25mm]{geometry}
\usepackage{fontspec}
\setmainfont{Latin Modern Roman}
\usepackage{amsmath,amssymb,mathtools,unicode-math}
\setmathfont{Latin Modern Math}
\usepackage{enumitem,needspace,xurl,hyperref}
\setlength{\parindent}{0pt}
\setlength{\parskip}{5pt}
\setlength{\emergencystretch}{3em}
\newcommand{\N}{\mathbb{N}}
\newcommand{\Z}{\mathbb{Z}}
\newcommand{\Q}{\mathbb{Q}}
\newcommand{\R}{\mathbb{R}}
\newcommand{\C}{\mathbb{C}}
\newcommand{\F}{\mathbb{F}}
\newcommand{\A}{\mathbb{A}}
\newcommand{\PP}{\mathbb{P}}
\newcommand{\E}{\mathbb{E}}
\newcommand{\eps}{\varepsilon}
\newcommand{\dd}{\,\mathrm{d}}
\begin{document}
\section*{UnsolvedMath ID: 2100707}
\textbf{Source problem code:} catalogue entry
\subsection{Definitions and mathematical statement}
Let \(T_{1,\hbar},\ldots,T_{n,\hbar}\) be a quantum integrable system with principal symbols \(f_1,\ldots,f_n\) on \((M,\omega)\), and set \(F=(f_1,\ldots,f_n)\). In a self-adjoint strongly commuting realization, write \(\Sigma_\hbar\subseteq\R^n\) for its joint spectrum. In a discrete-spectrum setting this is the set of tuples \((\lambda_1,\ldots,\lambda_n)\) for which a common nonzero eigenvector satisfies \(T_{j,\hbar}u=\lambda_ju\) for all \(j\). The semiclassical spectral data are the family \((\Sigma_\hbar)\) as \(\hbar\to0\), with the precision and multiplicity information to be specified.

Pelayo asks what information about the principal symbols, and in particular which symplectic invariants of \((M,\omega,F)\), can be recovered from those data. For a specified invariant \(I\), such as an integer \(z\) or matrix \(A\), the determination question is whether two admissible quantum systems with identical specified spectral data must satisfy
\[
I(M,\omega,F)=I(M',\omega',F').
\]
A reconstruction version seeks a rule \(R\) with
\[
R\bigl((\Sigma_\hbar)_{\hbar\to0}\bigr)=I(M,\omega,F).
\]
For example, an invariant under ordered symplectic equivalence is unchanged by a symplectomorphism \(\Phi\) satisfying \(F'\circ\Phi=F\); the source does not specify whether additional reparametrizations of the integrals are allowed.

No particular invariant, quantum calculus, geometric class, spectral precision, or meaning of effective computation is fixed, so this is a general inverse-spectral research question rather than a universal assertion that every invariant is recoverable. The supplied partial-progress summary reports recovery of semitoric invariants, including integer/matrix data, from semiclassical joint spectra, while saying the general problem remains open. It supplies neither the exact semitoric hypotheses nor a supporting theorem citation, so that claim is not extended to all semitoric or arbitrary integrable systems here.
\subsection{Short English statement}
Which classical symplectic invariants can be determined or computed from a quantum integrable system's semiclassical joint spectrum? The supplied summary reports semitoric progress but leaves the general recovery problem open.
\subsection{Sources}
\begin{itemize}
\item Local UnsolvedMath registry, permanent ID \(2100707\), exact source code \texttt{catalogue entry}, ``Open Problems in Integrable Systems --- Integrability and Quantisation.''
\item \emph{Open Problems, Questions, and Challenges in Finite-Dimensional Integrable Systems} (2018), Problem 7.7, attributed to \'{A}. Pelayo; \url{https://arxiv.org/abs/1804.03737}. Source-TeX extraction; recorded access 2026-07-29.
\item Supplied research summary refers to a semitoric quantum inverse spectral theorem but provides no complete authoritative citation or exact theorem hypotheses.
\end{itemize}
\end{document}
